\section{\it Decision Tree}

{\it Decision Tree} merupakan salah satu algoritma klasifikasi yang bekerja dengan memetakan data ke dalam kelas-kelas tertentu menggunakan urutan aturan keputusan berbentuk pola ``jika-maka'' (\textit{if-then}) (Tan, Steinbach, \& Kumar, 2019). Algoritma ini memilah kumpulan data (\textit{dataset}) menjadi kelompok-kelompok yang lebih spesifik dan seragam berdasarkan fitur (kolom) yang paling efektif dalam membedakan target kelas. Proses pemilahan ini dilakukan secara berulang (\textit{rekursif}) hingga memenuhi kriteria penghentian, misalnya saat seluruh data di dalam satu ujung cabang (\textit{node}) sudah memiliki label kelas yang sama, atau jumlah data terlalu sedikit untuk dipilah kembali (Han, Pei, \& Kamber, 2011).

Pada setiap percabangan, algoritma memilih fitur terbaik dengan mengukur tingkat "ketidakmurnian" atau ketidakpastian (\textit{impurity}) data, menggunakan metode seperti \textit{Information Gain}, \textit{Gini Index}, atau \textit{Gain Ratio}. Khusus untuk \textit{Information Gain}, rumus yang digunakan adalah sebagai berikut:
$$
IG(S, A) = Entropy(S) - \sum_{v \in Values(A)} \frac{|S_v|}{|S|} \cdot Entropy(S_v)
$$  
dengan $Entropy(S) = - \sum_{i=1}^{c} p_i \log_2 p_i$, di mana $p_i$ merupakan persentase proporsi data dalam suatu kelas. Entropi berfungsi untuk mengukur seberapa acak atau tidak pastinya sebuah himpunan data. Dengan demikian, \textit{Information Gain} pada dasarnya mengukur seberapa besar ketidakpastian tersebut berkurang setelah data dipilah menggunakan fitur $A$ (Quinlan, 1986).

Sebagai ilustrasi, simulasi perhitungan berikut diadaptasi dari \textit{dataset} klasik pada buku rujukan standar \textit{Machine Learning} (Mitchell, 1997) yang angka-angkanya dipetakan ke dalam konteks pengawasan penyewa. Misalkan terdapat 14 data penyewa ($S$) dengan fitur ``Lama Menginap'' dan ``Metode Pembayaran'', serta label kelas yang menandakan apakah penyewa tersebut tergolong normal (0) atau tidak normal (1). Mengingat pada studi kasus ini jumlah kelas tidak normal adalah minoritas, diasumsikan himpunan data terdiri atas 9 penyewa kelas normal (0) dan 5 penyewa kelas tidak normal (1). Pemilihan fitur mana yang paling cocok untuk membelah data pertama kali dilakukan dengan menghitung nilai \textit{Information Gain} (IG).

\begin{itemize}
    \item \textbf{Menghitung ketidakpastian awal (Entropi $S$)} \\
    Entropi awal dari keseluruhan himpunan 14 data penyewa dihitung sebagai berikut.
    $$
    Entropy(S) = - \left(\frac{9}{14} \log_2 \frac{9}{14}\right) - \left(\frac{5}{14} \log_2 \frac{5}{14}\right) \approx 0{,}940
    $$
    
    \item \textbf{Mengevaluasi fitur ``Lama Menginap''} \\
    Jika himpunan data dipecah menggunakan lama menginap, asumsikan data terbagi menjadi dua kelompok:
    \begin{itemize}
        \item \textbf{Kelompok 1--2 malam (7 penyewa)}, terdiri dari 3 Normal (0) dan 4 Tidak Normal (1). Karena perbandingannya cukup bervariasi, ketidakpastiannya dihitung sebagai berikut:
        $$
        Entropy(1\text{--}2\text{ malam}) = - \left(\frac{3}{7} \log_2 \frac{3}{7}\right) - \left(\frac{4}{7} \log_2 \frac{4}{7}\right) \approx 0{,}985
        $$
        \item \textbf{Kelompok $\geq$ 3 malam (7 penyewa)}, didominasi oleh kelas normal, yakni 6 Normal (0) dan 1 Tidak Normal (1). Ketidakpastiannya lebih rendah:
        $$
        Entropy(\geq3\text{ malam}) = - \left(\frac{6}{7} \log_2 \frac{6}{7}\right) - \left(\frac{1}{7} \log_2 \frac{1}{7}\right) \approx 0{,}592
        $$
    \end{itemize}
    
    \item \textbf{Menghitung \textit{Information Gain} (Lama Menginap)} \\
    Pengurangan ketidakpastian dihitung dengan mengurangkan entropi awal dengan rata-rata tertimbang (\textit{weighted average}) dari kelompok yang baru terbentuk:
    $$
    IG(S, Lama) = Entropy(S) - \left( \frac{|1\text{--}2\text{ malam}|}{|S|} \times Entropy(1\text{--}2\text{ malam}) \right) - \left( \frac{|\geq3\text{ malam}|}{|S|} \times Entropy(\geq3\text{ malam}) \right)
    $$
    $$
    IG(S, Lama) = 0{,}940 - \left( \frac{7}{14} \times 0{,}985 \right) - \left( \frac{7}{14} \times 0{,}592 \right) = 0{,}151
    $$
\end{itemize}

Setelah mendapatkan nilai evaluasi untuk Lama Menginap, perhitungan serupa dilakukan untuk fitur kandidat lainnya, yaitu ``Metode Pembayaran''. Asumsikan fitur ini memecah data menjadi kelompok ``Tunai'' dan ``Kartu''.

\begin{itemize}
    \item \textbf{Mengevaluasi fitur ``Metode Pembayaran''} \\
    Jika data dipecah berdasarkan metode pembayaran, distribusinya menjadi:
    \begin{itemize}
        \item \textbf{Kelompok Tunai (8 penyewa)}, terdiri dari 6 Normal (0) dan 2 Tidak Normal (1). Ketidakpastiannya adalah:
        $$
        Entropy(Tunai) = - \left(\frac{6}{8} \log_2 \frac{6}{8}\right) - \left(\frac{2}{8} \log_2 \frac{2}{8}\right) \approx 0{,}811
        $$
        
        \item \textbf{Kelompok Kartu (6 penyewa)}, terdiri dari 3 Normal (0) dan 3 Tidak Normal (1). Karena seimbang sempurna, ketidakpastiannya maksimal:
        $$
        Entropy(Kartu) = - \left(\frac{3}{6} \log_2 \frac{3}{6}\right) - \left(\frac{3}{6} \log_2 \frac{3}{6}\right) = 1
        $$
    \end{itemize}
    
    \item \textbf{Menghitung \textit{Information Gain} (Metode Pembayaran)} \\
    Sama seperti sebelumnya, pengurangan ketidakpastian untuk fitur ini adalah:
    $$
    IG(S, Metode) = Entropy(S) - \left( \frac{|\text{Tunai}|}{|S|} \times Entropy(Tunai) \right) - \left( \frac{|\text{Kartu}|}{|S|} \times Entropy(Kartu) \right)
    $$
    $$
    IG(S, Metode) = 0{,}940 - \left( \frac{8}{14} \times 0{,}811 \right) - \left( \frac{6}{14} \times 1 \right)
    $$
    $$
    IG(S, Metode) = 0{,}940 - 0{,}463 - 0{,}428 = 0{,}049
    $$
\end{itemize}

Berdasarkan perbandingan matematis di atas, pemisahan menggunakan fitur ``Lama Menginap'' berhasil mengurangi ketidakpastian sebesar 0,151, sedangkan ``Metode Pembayaran'' hanya 0,049 (Mitchell, 1997). Karena ``Lama Menginap'' memberikan pengurangan acak yang lebih baik ($0{,}151 > 0{,}049$), fitur tersebut secara otomatis dipilih oleh algoritma sebagai titik pembagi pertama (\textit{root node}) pada struktur pohon keputusan. Setelah titik pembagi awal ditetapkan, pohon akan terus mempercabangkan data. Penting untuk dicatat bahwa pada proses ini, algoritma tidak membuang kolom apa pun secara otomatis. Sistem akan memanfaatkan seluruh informasi relevan yang telah dikonversi (\textit{encoding}) dan mengevaluasi setiap kolom tersebut pada percabangan-percabangan selanjutnya hingga seluruh cabang bermuara pada kesimpulan akhir (daun/\textit{leaf}). Hasil akhir dari \textit{Decision Tree} adalah prediksi kelas untuk penyewa baru berdasarkan jalur kondisi yang telah dilewatinya. Semakin murni kondisi di ujung daun (misalnya berisi 100\% data tidak normal), semakin kuat keyakinan model terhadap prediksinya. Dalam penerapannya, model umumnya mengandalkan persentase kelas di ujung cabang sebagai nilai probabilitas. Sebagai contoh, jika suatu daun di ujung pohon memuat 8 data historis yang terdiri atas 6 kelas normal dan 2 kelas tidak normal, maka peluang probabilitas penyewa baru yang masuk ke daun tersebut untuk menjadi normal adalah 75\% (6/8). Angka probabilitas ini kemudian dapat dijadikan pijakan (\textit{threshold}) bagi pihak manajemen untuk menentukan tindakan; misalnya membunyikan alarm waspada jika peluang tidak normal menyentuh angka lebih dari 30\%.

Kelebihan utama dari algoritma \textit{Decision Tree} adalah kemampuannya memvisualisasikan logika pengambilan keputusan. Karena bentuknya menyerupai bagan alir dan cara kerjanya sangat identik dengan proses penalaran manusia, hasil dari model ini sangat mudah dipahami dan dijelaskan kepada pihak operasional maupun manajemen (Loh, 2011).